\chapter{Principal Component Analysis}\label{chap:pca}

\begin{tcolorbox}[colback=yellow!20,colframe=yellow!60!black,
  boxrule=0.4pt,arc=1mm,left=5pt,right=5pt,top=3pt,bottom=3pt,
  before skip=0pt,after skip=0pt]
\small\textbf{Start of lecture: PCA (Spinelli, Lecture 02)}
\end{tcolorbox}

\smallskip
So far, we have computed weights in a given basis. We now ask how to
choose useful directions from the data, using matrix transformations
and reconstruction error.

\section{Linear transformations}

A transformation takes an input vector $x$ and produces an output
vector $y$. We write it as $y=Ax$, where $A$ specifies how to combine
the input coordinates. If the input has $p$ coordinates and the output
has $q$, then $A$ has $q$ rows and $p$ columns.
Depending on its entries, this multiplication can stretch, rotate,
reflect, or project a point. Here $A$ denotes a general transformation
matrix; we will continue to use $C$ for a matrix of chosen directions.

\begin{frameddefn}{Linear transformation}\label{def:linear-transformation}
For a fixed matrix $A\in\mathbb{R}^{q\times p}$, the map
\[
  T:\mathbb{R}^p\to\mathbb{R}^q,\qquad T(x)=Ax
\]
is a linear transformation. The input $x$ has $p$ coordinates and the
output $y=Ax$ has $q$ coordinates. Each row of $A$ forms one output
coordinate as a weighted sum of the input coordinates.
For any input vectors $u,v$ and any real number $a$, it satisfies
\[
  T(u+v)=T(u)+T(v),\qquad T(au)=aT(u).
\]
These are the two properties meant by \emph{linear}: transforming a
sum or a scaled vector gives the same result as transforming its
parts first. In particular, the origin is mapped to the origin.
\end{frameddefn}

\Needspace{8\baselineskip}
\begin{framedexample}{Reading what a matrix does}
Suppose we want to double the first coordinate of $x=(2,1)^\top$
and halve its second coordinate. The first row of $A$ uses the
coefficients $2$ and $0$; the second uses $0$ and $0.5$.
Multiplying one row at a time gives
\[
  A=\begin{pmatrix}2&0\\0&0.5\end{pmatrix},\qquad
  y=Ax=\begin{pmatrix}2\cdot2+0\cdot1\\0\cdot2+0.5\cdot1\end{pmatrix}
      =\begin{pmatrix}4\\0.5\end{pmatrix}.
\]
The first row doubles the first coordinate; the second row halves the
second coordinate. The transformed vector is $y=(4,0.5)^\top$.
\end{framedexample}

\Needspace{10\baselineskip}
\subsection{What different matrices do}

Apply each of the following matrices to the same input $x=(2,1)^\top$.
The output makes the geometric effect concrete.

\begin{center}
\begin{tabularx}{\linewidth}{@{}lccX@{}}
\toprule
\textbf{Type} & $A$ & $y=Ax$ & \textbf{Effect} \\
\midrule
Scaling & $\begin{pmatrix}2&0\\0&0.5\end{pmatrix}$
  & $\begin{pmatrix}4\\0.5\end{pmatrix}$
  & Double the first coordinate; halve the second. \\[9pt]
Rotation & $\begin{pmatrix}0&-1\\1&0\end{pmatrix}$
  & $\begin{pmatrix}-1\\2\end{pmatrix}$
  & Rotate the point $90^\circ$ counterclockwise. \\[9pt]
Reflection & $\begin{pmatrix}-1&0\\0&1\end{pmatrix}$
  & $\begin{pmatrix}-2\\1\end{pmatrix}$
  & Reflect across the vertical axis. \\[9pt]
Shear & $\begin{pmatrix}1&1\\0&1\end{pmatrix}$
  & $\begin{pmatrix}3\\1\end{pmatrix}$
  & Add the second coordinate to the first. \\[9pt]
Projection & $\begin{pmatrix}1&0\\0&0\end{pmatrix}$
  & $\begin{pmatrix}2\\0\end{pmatrix}$
  & Keep the horizontal component. \\
\bottomrule
\end{tabularx}
\end{center}

For a square $p\times p$ matrix, $A^\top A=I_p$ means that its columns
are perpendicular unit vectors. It also means that multiplication by
$A$ preserves length, because
\[
  \lVert Ax\rVert_2^2
   =x^\top A^\top Ax=x^\top x=\lVert x\rVert_2^2.
\]
Applying the same calculation to the difference between two points
shows that their distance is preserved too. Rotations and reflections
in the table have this property. The scaling and shear matrices
change some lengths or angles.

Off-diagonal entries mix input coordinates. Both the rotation and the
shear above have such entries: their effect depends on the whole matrix.

\Needspace{16\baselineskip}
\subsection{How many directions survive?}

The projection matrix in the table maps every point $(a,b)^\top$
to $(a,0)^\top$: the second coordinate disappears. Its outputs
occupy a line rather than the whole plane. The \emph{rank} tells us
how many independent directions remain in the output.

\begin{frameddefn}{Rank of a matrix}\label{def:matrix-rank}
The rank of $A$ is the dimension of the space spanned by its columns:
the number of linearly independent output directions they provide.
We denote this number by $\operatorname{rank}(A)$.
\end{frameddefn}

For that projection, the first column is $(1,0)^\top$ and the second
is the zero vector. Only one direction remains, so its rank is $1$.
Both $(2,1)^\top$ and $(2,-1)^\top$ become $(2,0)^\top$.
The output alone cannot tell us which input we started with.

More generally, if $\operatorname{rank}(A)=K<p$, the outputs occupy
only $K$ independent directions, even if written with more than $K$
coordinates. Different inputs share an output, so no inverse can
recover every original input. Orthogonal projection is one such case.

\Needspace{10\baselineskip}
\section{Changing coordinates and projecting}

We now apply these transformations to representing data. Choose $K$
perpendicular unit directions for points with $p$ coordinates and put
them in the columns of $C$. For a point $x$, encoding computes the
weights $w=C^\top x$. Decoding combines the chosen directions with
those weights, giving $\hat{x}=Cw$.

\subsection{A complete basis keeps every direction}

First keep all $p$ directions of an orthonormal basis of
$\mathbb{R}^p$. Then $C$ is square, $C^\top$ is its inverse, and
$C^\top C=CC^\top=I_p$. Encoding followed by decoding gives
\[
  w=C^\top x,\qquad Cw=CC^\top x=x.
\]
The two multiplications undo each other. We still store $p$ numbers,
but in the new basis. Every point can be reconstructed exactly, and
distances between points are preserved.

\subsection{Fewer directions keep only a projection}

Now choose only $K<p$ of these directions. The matrix
$C\in\mathbb{R}^{p\times K}$ has fewer columns and their combinations
form a subspace $S$. The encoding and reconstruction are
\[
  C^\top C=I_K,\qquad
  w=C^\top x\in\mathbb{R}^K,\qquad
  \hat{x}=Cw=CC^\top x\in\mathbb{R}^p.
\]
The vector $w$ contains only $K$ numbers. The reconstructed vector
$\hat{x}$ has $p$ entries, but those entries come from only $K$
directions: $\hat{x}$ must lie in $S$. It equals $x$ when $x$ already
lies in $S$; otherwise it is the closest point in $S$.

\Needspace{19\baselineskip}
\begin{framedexample}{Two coordinates, one weight, two coordinates}
Suppose we want to describe $x=(2,3)^\top$ using just one weight.
Choose the horizontal unit direction $c_1=(1,0)^\top$ and set
$C=c_1$, a matrix with one column. Encoding gives the weight
\[
  w=C^\top x=1\cdot2+0\cdot3=2.
\]
Decoding uses that weight to produce
\[
  \hat{x}=Cw=2\begin{pmatrix}1\\0\end{pmatrix}
          =\begin{pmatrix}2\\0\end{pmatrix}.
\]
The process is
\[
  (2,3)^\top
    \quad\xrightarrow{\text{encode}}\quad 2
    \quad\xrightarrow{\text{decode}}\quad (2,0)^\top.
\]
The horizontal coordinate is kept. The vertical coordinate is lost:
even though the decoded vector has two entries, its second entry must
be zero for this choice of direction.
\end{framedexample}

\Needspace{10\baselineskip}
\section{Measuring reconstruction error}

We need a number that tells us how well a reduced representation
reconstructs the observations. For the point $x=(2,3)^\top$ and the
horizontal direction $C=(1,0)^\top$, decoding gives
$\hat{x}=(2,0)^\top$. Here $S$ is the horizontal axis. The vertical
segment in the diagram joins the reconstruction to the original point.

\begin{center}
\begin{tikzpicture}[font=\small,x=0.8cm,y=0.8cm]
  \draw[step=1,black!12,very thin] (0,0) grid (3,3);
  \draw[-{Stealth},NotesBlue,thick] (-0.4,0)--(3.8,0)
    node[right] {$S$};
  \draw[-{Stealth},black!50] (0,-0.3)--(0,3.6) node[above] {$x_2$};
  \draw[-{Stealth},black,thick] (0,0)--(2,3);
  \fill (2,3) circle (2pt) node[above right] {$x=(2,3)^\top$};
  \fill[NotesBlue] (2,0) circle (2pt)
    node[below] {$\hat{x}=(2,0)^\top$};
  \draw[-{Stealth},red!65!black,dashed,thick] (2,0)--(2,3)
    node[midway,right] {$r=(0,3)^\top$};
  \draw[black!60] (1.8,0)--(1.8,0.2)--(2,0.2);
  \node[below left] at (0,0) {$0$};
\end{tikzpicture}
\end{center}

The residual is $r=x-\hat{x}=(0,3)^\top$. Its length is $3$, the
perpendicular distance to $S$. The \emph{squared} reconstruction error
is $\lVert r\rVert_2^2=0^2+3^2=9$. For a plane in three dimensions,
the same idea measures the perpendicular segment from a point to the
plane.

\Needspace{13\baselineskip}
\begin{frameddefn}{Squared reconstruction error}\label{def:reconstruction-error}
In a dataset of $n$ observations with $p$ coordinates each,
$x^{(i)}$ is observation $i$ and $\hat{x}^{(i)}=Cw^{(i)}$ is its
reconstruction. Its squared reconstruction error is
\[
  E_i=\lVert x^{(i)}-\hat{x}^{(i)}\rVert_2^2
     =\sum_{j=1}^{p}\bigl(x_j^{(i)}-\hat{x}_j^{(i)}\bigr)^2.
\]
Here $j$ indexes a coordinate. Each difference
$x_j^{(i)}-\hat{x}_j^{(i)}$ is the error in that coordinate; $E_i$
adds their squares. It is zero when the observation is reconstructed
exactly.
\end{frameddefn}

\subsection{From one observation to the dataset}

Suppose the $K$ directions in $C$ are already chosen. We want a weight
vector for each of the $n$ observations, so we find the weights that
make the average of their squared errors as small as possible.
The symbol $\min$ below asks for this smallest average error:
\[
  \min_{w^{(1)},\ldots,w^{(n)}}
    \frac{1}{n}\sum_{i=1}^{n}
      \lVert Cw^{(i)}-x^{(i)}\rVert_2^2.
\]
When the columns of $C$ are perpendicular unit vectors, the best
weights are $w^{(i)}=C^\top x^{(i)}$. Each reconstruction is therefore
$Cw^{(i)}=C(C^\top x^{(i)})=CC^\top x^{(i)}$.
Substituting this expression into the error formula gives
\[
  \boxed{g(C)=\frac{1}{n}\sum_{i=1}^{n}
    \lVert CC^\top x^{(i)}-x^{(i)}\rVert_2^2.}
\]
We call this quantity $g(C)$ because it assigns one error value to
each choice of direction matrix $C$. The weights have already been
chosen as well as possible for those directions. To reduce the error
further, we must try different directions. Comparing their values of
$g(C)$ tells us which choice reconstructs the dataset better:
smaller $g(C)$ means smaller average error.

\Needspace{10\baselineskip}
\section{Learning the basis from the data}

So far, the directions in $C$ were given and we found the weights.
We now choose the directions themselves from the data. After deciding
how many weights $K$ to keep per observation, we compare choices of
$K$ perpendicular unit directions using $g(C)$. We want the directions
whose best reconstructions have the smallest average squared error.

For PCA, first mean-centre the data: subtract each feature's average
from that feature in every observation. Below, $x^{(i)}$ denotes a
centred observation. After reconstructing a centred point, add the
original mean vector to return to the original coordinates.

\Needspace{38\baselineskip}
\begin{framedexample}{Comparing two choices of direction}
Suppose we want to describe each of the following four observations
with only one weight:
\[
  (-3,-1)^\top,\quad (-3,1)^\top,\quad
  (3,-1)^\top,\quad (3,1)^\top.
\]
Their mean is $(0,0)^\top$, so they are already centred. We compare
two possible matrices, each containing one unit direction:
\[
  C_h=\begin{pmatrix}1\\0\end{pmatrix}
  \quad\text{(horizontal)},\qquad
  C_v=\begin{pmatrix}0\\1\end{pmatrix}
  \quad\text{(vertical)}.
\]
The black dots are the same data in both panels; blue dots show their
reconstructions. The dashed segments join each reconstruction to its
original point.

\begin{center}
\begin{tikzpicture}[font=\small]
  \begin{scope}[x=0.53cm,y=0.53cm]
    \node[font=\small\bfseries] at (0,2.5) {Keep the horizontal direction};
    \draw[-{Stealth},black!40] (0,-1.9)--(0,1.9);
    \draw[-{Stealth},NotesBlue,thick] (-4,0)--(4,0);
    \foreach \xx in {-3,3} {
      \draw[red!65!black,dashed] (\xx,-1)--(\xx,1);
      \fill[NotesBlue] (\xx,0) circle (2.5pt);
      \foreach \yy in {-1,1} {\fill (\xx,\yy) circle (2pt);}
    }
    \node[below right] at (0,0) {$0$};
  \end{scope}
  \begin{scope}[shift={(7.4cm,0)},x=0.53cm,y=0.53cm]
    \node[font=\small\bfseries] at (0,2.5) {Keep the vertical direction};
    \draw[-{Stealth},black!40] (-4,0)--(4,0);
    \draw[-{Stealth},NotesBlue,thick] (0,-1.9)--(0,1.9);
    \foreach \yy in {-1,1} {
      \draw[red!65!black,dashed] (-3,\yy)--(3,\yy);
      \fill[NotesBlue] (0,\yy) circle (2.5pt);
      \foreach \xx in {-3,3} {\fill (\xx,\yy) circle (2pt);}
    }
    \node[below right] at (0,0) {$0$};
  \end{scope}
\end{tikzpicture}
\end{center}

For the point $x=(3,1)^\top$, the horizontal choice gives weight
$C_h^\top x=3$ and reconstruction $(3,0)^\top$. The residual is
$(0,1)^\top$, so its squared error is $0^2+1^2=1$.
The vertical choice gives weight $C_v^\top x=1$ and reconstruction
$(0,1)^\top$. Its residual is $(3,0)^\top$, with squared error
$3^2+0^2=9$.

The other three points have the same squared errors: horizontal
projection removes a coordinate of magnitude $1$, whereas vertical
projection removes one of magnitude $3$. The average errors are
\[
  g(C_h)=\frac{1+1+1+1}{4}=1,\qquad
  g(C_v)=\frac{9+9+9+9}{4}=9.
\]

Both choices use one weight. The horizontal direction reconstructs
these observations better: it preserves their larger horizontal spread.
\end{framedexample}

\Needspace{20\baselineskip}
\begin{frameddefn}{Principal Component Analysis}\label{def:pca}
For mean-centred numerical data and a chosen dimension $K<p$,
Principal Component Analysis (PCA) learns $K$ perpendicular unit
directions that minimise the average squared reconstruction error.
Here $p$ is the original number of features, and $K$ is the number
of weights we keep per observation. The columns of $C$ contain the
learned directions. Each centred observation is encoded by
$w^{(i)}=C^\top x^{(i)}$ and reconstructed by $\hat{x}^{(i)}=Cw^{(i)}$.
\end{frameddefn}

PCA is \emph{unsupervised}: it learns the directions from the input
observations, without target labels. We specify how many directions to
keep; the data determine which directions are useful. Mean-centering
places the data cloud around the origin, so a subspace through the
origin can describe variation around the data's average.

The next question is how to find the directions that minimise $g(C)$.

\medskip
{\small\textit{Source: Spinelli, Lecture 02 (PCA), slides 2--8.}}
